{\large\textbf{QUESTION T1 [10 pts]. Three problems}}

This question consists of three independent problems.

\textbf{T1-A [3.0 pts]. Hydrostatic gate}

\begin{center}\includegraphics[width=0.49\linewidth]{figures/IPhO_2026_Q1_fig1.png}\end{center}

\begin{center}Fig. 1a.\end{center}

Two water reservoirs are separated by a vertical wall MN, as shown in the
figure 1a. The left reservoir can receive water from a source. It must be
ensured that the maximum possible difference between the water levels is
$\Delta h = 1.41\,\mathrm{m}$. To achieve this, a slot of vertical size
$a\sqrt{2}/2$ is cut in MN and sealed with a solid cubic block of side length
$a$. The block has density $3\rho_0$, with $\rho_0$ the density of water. The
block is fully submerged, fixed to the wall at $O$ and can rotate without
friction about the axis perpendicular to the plane of the figure that passes
through point $O$.

\textbf{T1-A1 [3.0 pts.]} Calculate the value of $a$ that ensures the maximum
permissible difference in water levels is not exceeded.

\textbf{T1-B [3.5 pts]. Electron-positron pair}

\begin{center}\includegraphics[width=0.20\linewidth]{figures/IPhO_2026_Q1_fig2.png}\end{center}

\begin{center}\texttt{Fig. 1b.}\end{center}

At a certain instant of time, a positron $e^+$ is located at a distance of
$100a_0$ from an electron $e^-$. The particles are moving in such a way that
their velocities are antiparallel at that instant, perpendicular to their
separation, as shown in figure 1b. It is known that each particle has an
angular momentum of magnitude $\mu\hbar$ with respect to the system's center of
mass, where $\mu$ is a dimensionless numerical factor. Assume that the only
interaction between $e^+$ and $e^-$ is electrostatic (both particles have the
same mass $m$ and equal magnitude of the charge but of opposite sign).
Furthermore, assume that this system remains isolated and that its dynamics are
classical and non-relativistic.

\textit{The Bohr radius $a_0$ is given by $a_0 = 4\pi\epsilon_0\hbar^2/(me^2)$.}

\textbf{T1-B1 [1.0 pts.]} If $\mu = 4$, the system is bound, meaning that the
particles move in a closed orbit around the system's center of mass. Find the
maximum separation distance between $e^+$ and $e^-$ in terms of $a_0$.

\textbf{T1-B2 [2.5 pts].} If $\mu = \frac{15}{2}$, the system is unbound,
meaing that the particles are no longer in closed orbit around the system's
center of mass. Let $\vec{u}_\infty$ be the velocity of $e^+$ relative to $e^-$
as the separation distance between them tends to infinity. Calculate the angle
between $\vec{u}_\infty$ and the initial line of motion of $e^+$. Give your
answer in degrees.

\textit{Hint 1: The eccentricty $\varepsilon$ of a conic trajectory for the
present case is given by}
\[
\varepsilon = \sqrt{1 + \frac{4L^2E}{k^2e^4m}}
\]
\textit{where E and L are the total energy and the magnitude of the total
angular momentum, respectively. k is Coulombs constant $k = 1/(4\pi\epsilon_0)$.}

\textit{Hint 2: The equation of a conic in polar coordinates is}
\[
r = \frac{a}{1 - \varepsilon\cos(\theta)}
\]

\textbf{T1-C [3.5 pts]. Photodissociation of ozone}

\begin{center}\includegraphics[width=0.37\linewidth]{figures/IPhO_2026_Q1_fig3.png}\end{center}

\begin{center}Fig. 1c.\end{center}

A photon with angular frequency $\omega$ strikes an ozone molecule O$_3$ at
rest, dissociating it (breaking it apart) into an oxygen molecule O$_2$ and an
oxygen atom O, as shown in figure 1c. The photon is absorbed in this process.
Let $U_i$ and $U_f$ be the ground state energies of the ozone and oxygen
molecules, respectively. Assume that this system remains isolated and that its
dynamics is classical (potential energies do not contribute to mass and the
motion of the Oxygen atoms are non-relativistic). Use the relation $p = E/c$
for the linear momentum of the photon, where $E$ is its energy.

\textbf{T1-C1 [2.5 pts].} If the angle that the momentum of the outgoing O$_2$
with respect to the incident photon is $\theta$, determine the minimum angular
frequency $\omega_{\mathrm{min}}$ required for this dissociation to occur at
this angle. Write your answer in terms of $\hbar$, $c$, $\theta$,
$\Delta U = U_f - U_i$, and the mass $m$ of an oxygen atom.

\textbf{T1-C2 [1.0 pts].} If $\theta = \pi/6$, calculate the value of
$\hbar\omega_{\mathrm{min}} - \Delta U$. Give your answer in electronvolts
(eV). Take $\Delta U = 1.10\,e\mathrm{V}$ and $m = 16.0\,\mathrm{amu}$.

\textit{Hint:}
\[
(1+x)^r = 1 + rx + \frac{r(r-1)}{2!}x^2 + \frac{r(r-1)(r-2)}{3!}x^3 + \cdots
\]
